\section{Metodologia}

La metodologia combina modelado, simulacion, ablacion y validacion estadistica. El
procedimiento se organiza por modulos, de forma que cada paso de la escalera M0--M5 tenga
una pregunta falsable, una ecuacion minima, un escenario y una lectura de resultados. El
flujo general se resume en cinco pasos.

\begin{enumerate}[leftmargin=2em]
  \item Definir escenarios reproducibles con semillas fijas, cargas, robots, obstaculos,
  conectividad y demanda de coalicion.
  \item Ejecutar tratamientos comparables: greedy, centralizado limitado, Smith,
  Smith-QR, proxies clasicos y baselines MARL.
  \item Registrar trayectorias, asignaciones, preferencias, precios, estados de carga,
  esfuerzo, distancia, contactos y tiempos de completacion.
  \item Agregar metricas por escenario y semilla, usando comparaciones pareadas cuando
  los metodos comparten condiciones iniciales.
  \item Ejecutar una suite de hipotesis que produce CSV, JSON, tablas LaTeX y una figura
  de resumen.
\end{enumerate}

La configuracion experimental queda fuera del codigo de control. Los escenarios
reproducibles se describen en \texttt{experiments/*/config.yaml}; los parametros
transversales y modelos aprendidos se documentan en \texttt{configs/}; y las campanas
H10--H12 usan scripts de orquestacion con salidas cerradas en \texttt{results/campaigns/}.
Esta separacion permite cambiar semillas, radio de comunicacion, cargas, fallos o
tratamientos sin modificar asignadores ni controladores.

El motor principal es Python porque permite miles de corridas rapidas y deterministas.
CoppeliaSim queda como capa de plausibilidad robotica e integracion visual. Cuando una
campana usa render sintetico de respaldo, se reporta como evidencia de artefactos y no
como simulacion fisica. Esta separacion evita que el TFM dependa de una simulacion pesada
para probar cada hipotesis, pero conserva una ruta de integracion con un entorno robotico
mas rico.

\subsection*{Protocolo incremental por fases}

Para cada modulo se ejecuta el mismo patron experimental. Primero se fija el caso minimo:
numero de robots, cargas, grafo, obstaculos, demanda y tratamiento. Segundo se ejecuta el
metodo base del nivel anterior. Tercero se activa el nuevo mecanismo: quorum entero,
capacidad efectiva, rol fisico, campo unico o bateria. Cuarto se comparan metricas
pareadas, evitando declarar dominancia cuando el intervalo de confianza no lo soporta.

\begin{table}[H]
\centering
\footnotesize
\begin{tabularx}{\textwidth}{@{}p{0.12\textwidth}p{0.25\textwidth}Y Y@{}}
\toprule
Modulo & Escenario minimo & Metrica principal & Salida en el reporte \\
\midrule
M0 & Greedy/centralizado/Smith bajo condiciones nominales y R3. & Captura, entregas, distancia desperdiciada. & Tabla comparativa y fallo declarado. \\
M1 & Varias cargas con demanda de quorum. & Deficit entero y tareas completadas. & Water-filling y clearing. \\
M1.3 & Cargas ligeras y pesadas con robots iguales. & Priorizacion, throughput y deadlines. & Barridos de carga/escasez. \\
M1.4 & Carga que requiere fuerza y torque. & Residual wrench y saturacion. & Gate algebraico y H11. \\
M2--M3 & Carga rectangular con slots. & Error de formacion y contacto valido. & Escenas, roles derivados y H12. \\
M4 & Obstaculos, fallos y comunicacion local. & Recuperacion, congestion y switches. & H10--H12. \\
M5 & Robot con bateria baja. & Margen de retornabilidad. & Gate G7 y extension futura. \\
\bottomrule
\end{tabularx}
\caption{Contrato experimental incremental por modulo.}
\label{tab:protocolo-modular}
\end{table}
